\chapter{Część badawcza}
    \section{Wybór otwartoźródłowego modelu językowego}
        \subsection{Przegląd dostępnych modeli (np. Llama, Mistral, Falcon, itp.)}
        \subsection{Kryteria wyboru modelu (wydajność, rozmiar, licencja, możliwości fine-tuning'u)}
        \subsection{Uzasadnienie wyboru konkretnego modelu}
    \section{Metodologia badawcza}
        \subsection{Opis wykorzystanych zbiorów danych}
            \subsubsection{Zbiór treningowy}
            \subsubsection{Zbiór walidacyjny}
            \subsubsection{Zbiór testowy}
        \subsection{Przygotowanie danych (Preprocessing)}
        \subsection{Proces fine-tuning'u modelu}
            \subsubsection{Architektura zadania klasyfikacyjnego}
            \subsubsection{Parametry i hiperparametry uczenia}
            \subsubsection{Wykorzystane narzędzia i biblioteki}
        \subsection{Metryki oceny skuteczności}
            \subsubsection{Dokładność (Accuracy)}
            \subsubsection{Precyzja (Precision)}
            \subsubsection{Pełność (Recall)}
            \subsubsection{Miara F1 (F1-Score)}
            \subsubsection{Inne istotne metryki (np. AUC)}
    \section{Implementacja i przeprowadzenie eksperymentów}
        \subsection{Środowisko eksperymentalne}
        Witajcie co tam  HALo!
        \subsection{Przebieg procesu fine-tuning'u}
        KUUURWA
    \section{Analiza skuteczności dostrojonego modelu}
        \subsection{Wyniki na zbiorze testowym}
        \subsection{Porównanie z wynikami istniejących metod (z Rozdziału 2)}
        \subsection{Analiza błędów klasyfikacji (False Positives, False Negatives)}
        \subsection{Ocena elastyczności modelu i odporności na nowe typy spamu} % Odporność na ewolucję zagrożeń
    \section{Analiza możliwości wdrożenia produkcyjnego}
        \subsection{Wymagania sprzętowe i czas wnioskowania (inference time)}
        \subsection{Koszty obliczeniowe (fine-tuning i wnioskowanie)}
        \subsection{Skalowalność rozwiązania}
        \subsection{Potencjalne wyzwania i ograniczenia wdrożeniowe}
        \subsection{Praktyczność zastosowania w porównaniu do tradycyjnych filtrów}
